\chapter*{INTRODUCTION}
\addcontentsline{toc}{chapter}{INTRODUCTION}

Cryptocurrencies were originally developed to make digital money transfers possible without the involvement of banks or central authorities. The main idea was to create a system in which users could send and receive value directly, without relying on a trusted third party to process or verify transactions \cite{nakamoto2008bitcoin}. This move toward decentralization replaced traditional institutional trust with a system in which transactions are governed by code and recorded on a shared digital network.

Today, however, the cryptocurrency market has developed into a much more complex system. While many projects still aim to build useful financial infrastructure, a large part of the market is shaped by social media attention and speculative trading \cite{mirtaheri2021identifying, hamrick2021examination, lamorgia2023doge}. New assets are often launched and traded globally before their actual use can be properly assessed. Prices can move sharply within minutes, trading activity is spread across many exchanges, and market conditions often differ strongly between large coins and smaller tokens \cite{makarov2020trading, liu2021risks}. For some traders, this volatility creates opportunities, but it also opens the door to organized manipulation that targets less experienced investors.

This thesis examines one specific form of manipulation: pump-and-dump schemes coordinated through messaging platforms. In such schemes, organizers build attention before an event and then announce a target asset to a group of traders at the same time. The signal is designed to create immediate buying pressure. If enough participants react quickly, the price may rise for a short period, allowing early buyers or insiders to sell before the market reverses. Previous studies show that these schemes are visible in cryptocurrency markets and are often organized through Telegram, Discord, and similar platforms \cite{kamps2018pump, xu2019anatomy, hamrick2021examination, li2021cryptocurrency, lamorgia2023doge}.

The research problem is not only whether pump-and-dump events exist. That question has already been studied in the literature. The more focused problem is whether the style and structure of a pump signal help explain the immediate market response after the coin is revealed. A signal can be short or detailed, calm or urgent, isolated or preceded by repeated hype messages. These differences matter because the message is not a normal piece of financial news. It is a coordination tool that tries to make many traders act at the same moment.

This creates a clear research gap. Existing studies provide strong evidence that cryptocurrency pump-and-dump schemes can be detected, measured, and linked to abnormal price and volume changes \cite{kamps2018pump, victor2019cryptocurrency, li2021cryptocurrency}. Other work studies the anatomy of these groups, the role of Telegram and Discord, and the possibility of predicting target coins or detecting manipulation in real time \cite{xu2019anatomy, lamorgia2020pump, hu2023sequence}. These contributions are important, but they do not fully explain whether the wording and timing of the public signal itself help drive the size of the immediate price reaction. This thesis therefore moves from identifying pump events to studying how their communication design is connected to market outcomes.

This problem is important for both market analysis and investor protection. Pump-and-dump schemes can create artificial price movements that do not reflect changes in the real value of the asset. The risk is especially high in thinner markets, where limited trading depth means that even modest buying pressure can move prices strongly. Market microstructure research shows that illiquidity increases price impact, so the same signal may produce very different reactions depending on the asset being targeted \cite{amihud2002illiquidity, kyle1985continuous}. For this reason, the thesis pays particular attention to the role of \textit{Market Capitalization}, which is used as a practical proxy for the market's ability to absorb coordinated demand.

The study is built around three hypotheses. First, more urgent pump signals are expected to have a positive estimated effect on short-term cryptocurrency returns after the coin reveal, especially in lower-cap market segments where prices are easier to move. Second, signals that create urgency and coordination are expected to have stronger estimated effects in lower-cap assets, especially small-cap assets, and these effects are expected to become weaker as \textit{Market Capitalization} increases. Third, observable market patterns before the event, combined with signal characteristics, are expected to contain information that can help predict \textit{Success}.

The empirical analysis uses a dataset that combines Telegram messages, high-frequency exchange data, and asset metadata. Historical Telegram messages are used to identify pump sessions and extract signal characteristics. One-minute OHLCV (Open High Low Close Volume) data from cryptocurrency exchanges is used to measure the market reaction around the reveal time, denoted as $T_0$. Historical CoinMarketCap metadata supplies additional information about the targeted assets. After validation and complete-case restrictions, the final causal sample contains 2,917 pump events.

The main outcome is the \textit{10-Minute Logarithmic Return} after the coin reveal. This short window is chosen because pump-and-dump schemes operate very quickly and because longer windows would include more unrelated market movements. The causal part of the analysis uses DML (Double Machine Learning) to estimate how signal characteristics are related to short-term returns after controlling for pre-reveal market conditions. The study also examines whether these effects differ across \textit{Market Capitalization} strata. A separate predictive part uses tree-based machine learning models, especially XGBoost (Extreme Gradient Boosting), to test whether \textit{Success}, defined as a positive \textit{10-Minute Logarithmic Return}, can be predicted up to the reveal.

The empirical work was carried out in Python using external libraries. Telegram data extraction used Telethon, exchange data extraction used CCXT, causal estimation used the econml library, and predictive modeling used tools such as XGBoost, LightGBM, Random Forests, and Optuna. These tools made it possible to combine unstructured message data, high-frequency market data, causal estimation, and prediction within one empirical pipeline.

The expected contribution is therefore both substantive and methodological. Substantively, the thesis treats pump messages as part of the manipulation mechanism, not only as labels that help locate suspicious market events. Methodologically, it separates causal explanation from forecasting. This distinction is important because a variable can help predict \textit{Success} without being the direct cause of the return. Keeping these tasks separate makes the interpretation of the empirical results clearer.

The following chapters present the research in detail. Chapter I reviews the literature on market manipulation, cryptocurrency pump-and-dump schemes, coordinated trading signals, market microstructure, and causal measurement. Chapter II describes the data collection process, event construction, and feature engineering. Chapter III explains the methodology, including the causal framework, heterogeneous-effect analysis, and predictive modeling. Chapter IV reports the empirical results, while Chapter V interprets the findings in relation to the hypotheses. The Conclusions summarize the main contributions and limitations, and the Appendix discusses the synthetic-control exercise considered during the design stage.
